\documentclass[aspectratio=169,11pt]{ctexbeamer}

\usetheme{Madrid}
\usecolortheme{default}
\setbeamertemplate{navigation symbols}{}
\setbeamertemplate{caption}[numbered]

\usepackage{graphicx}
\usepackage{booktabs}
\usepackage{array}
\usepackage{tabularx}
\usepackage{hyperref}

\title[水彩风景 LoRA 微调]{Stable Diffusion v1.5 水彩风景 LoRA 微调}
\subtitle{课程大作业中期汇报}
\author{刘家宁，熊禹翔，周航，朱昱醒，车昕宇}
\date{2026 年 5 月 6 日}

\newcommand{\img}[2]{\includegraphics[width=#1,keepaspectratio]{#2}}

\begin{document}

\begin{frame}
  \titlepage
\end{frame}

\begin{frame}{选题与研究目标}
  \begin{columns}[T,onlytextwidth]
    \begin{column}{0.58\textwidth}
      \begin{itemize}
        \item 选题：使用 LoRA 微调 Stable Diffusion v1.5，生成稳定的水彩风景插画。
        \item 目标风格：山湖、雪景、小镇、日落、花草等。
        \item 核心问题：基础模型只靠 \texttt{watercolor style} prompt 时，风格不稳定。
        \item 目标输出：提升水彩风格一致性，同时保留 prompt 控制能力和生成多样性。
      \end{itemize}
    \end{column}
    \begin{column}{0.38\textwidth}
      \centering
      \img{0.95\linewidth}{assets/beamer/examples/overall.jpg}

      \vspace{0.3em}
    \end{column}
  \end{columns}
\end{frame}

\begin{frame}{为什么需要微调：Base 生成不稳定（一）}
  \centering
  \begin{columns}[T,onlytextwidth]
    \begin{column}{0.32\textwidth}
      \centering
      \img{0.95\linewidth}{assets/beamer/base/01_seed-42_watercolor-style-a-small-boat-on-a-mountain-lake.png}

      {\scriptsize \texttt{a small boat on a mountain lake}}
    \end{column}
    \begin{column}{0.32\textwidth}
      \centering
      \img{0.95\linewidth}{assets/beamer/base/02_seed-42_watercolor-style-a-snowy-forest-path.png}

      {\scriptsize \texttt{a snowy forest path}}
    \end{column}
    \begin{column}{0.32\textwidth}
      \centering
      \img{0.95\linewidth}{assets/beamer/base/03_seed-42_watercolor-style-a-quiet-village-with-a-stone-bridge.png}

      {\scriptsize \texttt{a quiet village with a stone bridge}}
    \end{column}
  \end{columns}

  \vspace{0.6em}
  \begin{itemize}
    \item SD v1.5 Base 能理解小船、湖泊、桥、村庄等主体。
    \item 但风格稳定性不一致：有些图像偏摄影写实，有些才有明显水彩感。
  \end{itemize}
\end{frame}

\begin{frame}{为什么需要微调：Base 生成不稳定（二）}
  \centering
  \begin{columns}[T,onlytextwidth]
    \begin{column}{0.48\textwidth}
      \centering
      \img{0.72\linewidth}{assets/beamer/base/04_seed-42_watercolor-style-pine-trees-beside-a-lake-at-sunset.png}

      {\scriptsize \texttt{pine trees beside a lake at sunset}}
    \end{column}
    \begin{column}{0.48\textwidth}
      \centering
      \img{0.72\linewidth}{assets/beamer/base/05_seed-42_watercolor-style-wildflowers-beside-a-peaceful-river.png}

      {\scriptsize \texttt{wildflowers beside a peaceful river}}
    \end{column}
  \end{columns}

  \vspace{0.8em}
  \begin{itemize}
    \item 仅依赖 prompt 时，水彩风格是“可能出现”，不是“稳定出现”。
    \item LoRA 微调的目标：在保持主体控制能力的同时，提高水彩风格一致性。
  \end{itemize}
\end{frame}

\begin{frame}{为什么选择 LoRA}
  \small
  \begin{tabularx}{\textwidth}{l X}
    \toprule
    方案 & 取舍原因 \\
    \midrule
    LoRA & 只训练少量低秩适配参数，训练轻量、checkpoint 小，适合 100 张级别的小数据集风格微调。 \\
    全量微调 & 需要更新整个 SD v1.5，显存和时间成本更高；小数据下更容易过拟合，也更可能破坏原模型的 prompt 控制能力。 \\
    DreamBooth & 更适合学习特定主体或身份；本项目目标是水彩风景风格，不是绑定某一个具体对象。 \\
    \bottomrule
  \end{tabularx}

  \vspace{0.8em}
  \begin{block}{选择原则}
    本项目需要训练 Group90 和 A100--E100 多个模型，LoRA 更方便快速训练、保存和横向比较。
  \end{block}
\end{frame}

\begin{frame}{数据集结构：公共 90 张 + 个人 10 张}
  \begin{columns}[T,onlytextwidth]
    \begin{column}{0.48\textwidth}
      \textbf{公共数据集：90 张}
      \begin{itemize}
        \item A 类：山湖，18 张
        \item B 类：雪景，18 张
        \item C 类：小镇，18 张
        \item D 类：日落，18 张
        \item E 类：花草，18 张
      \end{itemize}
    \end{column}
    \begin{column}{0.48\textwidth}
      \textbf{每位成员：对应风格额外 10 张}
      \begin{itemize}
        \item 每人训练集：公共 90 张 + 自己负责风格的 10 张。
        \item 全组实际收集：90 + 5 $\times$ 10 = 140 张。
        \item 目的：公共数据保持均衡，个人数据体现对应风格增强。
      \end{itemize}
    \end{column}
  \end{columns}
\end{frame}

\begin{frame}{数据集示例（一）}
  \centering
  \begin{tabular}{ccc}
    \img{0.30\linewidth}{assets/beamer/manual/style_a.png} &
    \img{0.30\linewidth}{assets/beamer/manual/style_b.png} &
    \img{0.30\linewidth}{assets/beamer/manual/style_c.png} \\
    {\scriptsize 风格 1：山湖} &
    {\scriptsize 风格 2：雪景} &
    {\scriptsize 风格 3：小镇}
  \end{tabular}

  \vspace{0.8em}
  \begin{itemize}
    \item 数据集按 5 个水彩风景方向收集，每类 18 张公共图像。
    \item 每位成员在对应方向上额外补充 10 张，用于观察局部风格增强。
  \end{itemize}
\end{frame}

\begin{frame}{数据集示例（二）}
  \centering
  \begin{tabular}{cc}
    \img{0.36\linewidth}{assets/beamer/manual/style_d.png} &
    \img{0.36\linewidth}{assets/beamer/manual/style_e.png} \\
    {\scriptsize 风格 4：日落} &
    {\scriptsize 风格 5：花草}
  \end{tabular}

  \vspace{0.8em}
  \begin{itemize}
    \item 所有样例图统一处理为 512 $\times$ 512，保证后续训练尺寸一致。
    \item 正式训练前会继续检查水印、文字、主体裁切和风格混入问题。
  \end{itemize}
\end{frame}

\begin{frame}{统一标注与数据处理}
  \begin{columns}[T,onlytextwidth]
    \begin{column}{0.50\textwidth}
      \textbf{图像处理}
      \begin{itemize}
        \item 筛选标准：水彩风景主风格，避免摄影、厚涂、动漫、3D 渲染混入。
        \item 裁剪原则：保留主体、构图、纸张纹理，统一到 512 $\times$ 512。
        \item 质量检查：剔除水印、文字、低清晰度、主体严重裁切图片。
      \end{itemize}
    \end{column}
    \begin{column}{0.46\textwidth}
      \textbf{标注格式}
      \begin{itemize}
        \item 文件名：\texttt{public\_001.png}，\texttt{member\_a\_001.png}。
        \item Caption：统一风格词 + 简短主体描述。
        \item 来源和授权信息单独留存，方便报告引用和复查。
      \end{itemize}
    \end{column}
  \end{columns}

  \vspace{0.8em}
  \textbf{统一 caption 格式：}

  \vspace{0.3em}
  \texttt{watercolor landscape style, [subject description]}

  \vspace{0.4em}
  {\scriptsize 示例：\texttt{watercolor landscape style, a small boat on a mountain lake}}
\end{frame}

\begin{frame}{训练方案}
  \begin{columns}[T,onlytextwidth]
    \begin{column}{0.48\textwidth}
      \textbf{模型与方法}
      \begin{itemize}
        \item Base：Stable Diffusion v1.5
        \item 方法：LoRA 风格微调
        \item 分辨率：512 $\times$ 512
        \item \texttt{batch\_size}：4
        \item 资源规划：单张 RTX 4090
      \end{itemize}
    \end{column}
    \begin{column}{0.48\textwidth}
      \textbf{实验模型}
      \begin{itemize}
        \item Group90：公共 90 张，5 类各 18 张
        \item A100--E100：公共 90 + 对应风格额外 10 张
        \item 保存 checkpoint：20/40/60/80 epochs
      \end{itemize}
    \end{column}
  \end{columns}
\end{frame}

\begin{frame}{评估指标与个人影响分析}
  \small
  \begin{tabularx}{\textwidth}{l X}
    \toprule
    指标 & 计算方式与作用 \\
    \midrule
    CLIP Score & 用 CLIP 提取文本和生成图特征，计算余弦相似度；衡量 prompt 遵循能力 \\
    Style Similarity & 提取训练集图像特征并求中心，计算生成图到训练集中心的相似度；衡量水彩风格学习效果 \\
    Diversity & 计算生成图两两之间的 CLIP 特征距离或 LPIPS 距离；衡量是否模式坍缩 \\
    Nearest Neighbor & 对每张生成图在训练集中找最近邻并可视化；检查是否过拟合或复制训练图 \\
    Personal Influence & 比较个人模型和 Group90 到该成员额外 10 张数据中心的相似度差值；衡量对应风格增强 \\
    \bottomrule
  \end{tabularx}

\end{frame}

\begin{frame}{组内分工与实施计划}
  \small
  \begin{tabularx}{\textwidth}{l X}
    \toprule
    阶段 & 工作内容 \\
    \midrule
    数据收集 & 公共 90 张按 5 类各 18 张收集；每人负责对应风格额外 10 张 \\
    数据处理 & 筛图裁剪；统一整理 caption、来源授权与质量抽查 \\
    训练实验 & 先训练 Group90，再训练 A100--E100，并保存多个 checkpoint \\
    对比生成 & 使用 5 个固定个人主题 prompt，对 Base、Group90、个人模型做对比 \\
    指标评估 & 计算 CLIP Score、Style Similarity、Diversity、Nearest Neighbor、Personal Influence \\
    报告撰写 & 汇总图像对比、指标表格、过拟合分析和个人数据影响分析 \\
    \bottomrule
  \end{tabularx}
\end{frame}
